\chapter{Extensions to Deffuant Model}\label{app:extension}
There are various extensions to this model presented in this model that will not be covered but are worth mentioning. One such possible extension is known as the heterogeneous Deffuant–Weisbuch model \cite{chen2020convergence}. This model does not make the assumption that the r value is constant (or the mean is constant in the case of the noisy Deffuant model), instead each r value is unique between agents (hence the term heterogeneous). Traits such as convergence and a bounded mean-square error are seen in this version of the Deffuant model. Similar could also be done with the value of $\mu$, but seeing as the value of $\mu$ had little effect on the long term behaviour of the lattice in this report one would expect little to change between the findings of this report and this extension (although this is not confirmed).

Another extension is to assume that agents who do not communicate (that is, sites that are outside the specified r value) instead repel each other away from each other, rather than doing nothing \cite{lanchier2023deffuantopiniondynamicsattraction}. This is referred to as repulsion and is a good fit for the Deffuant model as it only really had attraction before. Mixes of repulsion and attraction are also an interesting extension of the Deffuant model.

Another possible extension is the implementation of zealots. These are lattice sites that would not change value regardless of communication with other sites. One would expect clusters to form around these sites and has some interesting consequences for other models and different dimensions, such as breaking conservation of magnetisation for dimension less than or equal to 2 \cite{castellano2009statistical}.

The models in this report can be used on other types of graphs and are not limited to square/rectangular, periodic lattices. Any sort of network or odd-shaped lattice could also be studied. Generation of random graphs with random weightings in between nodes may provide a more realistic scenario than a square lattice.

There are many other extensions to the Deffuant model out there, from hybrid models that combine existing models together, to new takes on old models. These are just some modifications that have been investigated or considered.